\section{Dashboard Design and Usability}
The dashboard (\texttt{dashboard/}) uses only D3.js v7 and the TopoJSON client, both vendored with
subresource-integrity hashes. It loads only the aggregated JSON exports, which must stay under 10\,MB: the
export keeps the most recent weeks that fit. Five views are linked: a district choropleth of incidence per
100,000 with alert markers, a weekly time series with a brush, forecast and bands, a seasonal heatmap, an
allocation view with a what-if panel, and a ranked early-warning table. Suppressed cells are hatched and every
view shows the data-freshness time.
\textbf{Design rationale.} Incidence uses the viridis scale, which stays ordered in greyscale and for the
common colour-vision deficiencies; alerts use the Okabe--Ito vermillion, which stands out against it. Counts
are shown per 100,000 on the map, so a large district does not look worse only for being large, but as raw
cases on the time series, which is what a store-keeper plans with. Suppressed cells are hatched rather than
left blank, so ``hidden'' never reads as ``zero''. Every forecast carries its 80\% band. District by week is
the finest level shown: it is what an officer acts on and the coarsest level that still locates an outbreak.
Below 560\,px the views stack in one column for a phone.

\textbf{District view (PRV-3).} Opening the dashboard with \texttt{?district=} and a district code loads that
district's bundle: its own series, alerts, forecast and allocation in full, and the rest of the country drawn
lighter as division totals.

% PERSONA: simulated persona walkthroughs (step 4)
